% ============================================================
% CHAPTER 7: CONCLUSION AND FUTURE WORK
% Publication-Grade LaTeX
% ============================================================

\chapter{Conclusion and Future Work}
\label{chap:conclusion}

% ============================================================
% 7.1 Summary of Contributions
% ============================================================
\section{Summary of Contributions}

This thesis has presented a comprehensive investigation into Named Entity Recognition for Sanskrit, addressing the fundamental challenge of limited annotated resources for this morphologically rich, low-resource language. Through systematic experimentation spanning 15+ fine-tuning and validation runs, we have developed a hybrid neuro-symbolic approach that achieves state-of-the-art performance while maintaining robust cross-domain generalisation.

\subsection{Primary Contributions}

\begin{enumerate}
    \item \textbf{First Correct DCS Sembank NER Extraction Methodology (to our knowledge)}: We developed and validated a novel extraction pipeline that correctly distinguishes entity mentions (col1 IDs) from entity references (col0 IDs) in the DCS Sembank \cite{hellwig2010dcs, hellwig2025sembank}, yielding 12,955 high-quality silver training examples from 120 diverse Sanskrit texts. This methodology resolves the common noun contamination that plagued previous attempts and represents a methodological contribution enabling future researchers to leverage the Sembank for entity-centric tasks.

    \item \textbf{Hybrid Training Strategy with Linguistic Regularisation}: We demonstrated that incorporating 15\% linguistic regularisation data (segmentation and lemmatisation tasks) during NER fine-tuning prevents catastrophic forgetting \cite{mccloskey1989catastrophic, goodfellow2013empirical} of base model capabilities while maintaining competitive in-domain performance. This strategy enables the Final NER model to achieve the best balance between in-domain accuracy (F1 = 0.852) and cross-domain generalisation (73.4\% PROPN recall), outperforming both pure NER training and large-scale multi-task learning \cite{caruana1997multitask}.

    \item \textbf{Comprehensive Evaluation Framework}: We established the first comprehensive NER baselines for Sanskrit, evaluating 15 models across in-domain (Mah\={a}n\={a}ma \cite{sarkar2025mahanama} test set) and cross-domain (Pa\~{n}catantra \cite{zeman2023ud}) settings, with statistical significance analysis using bootstrap confidence intervals \cite{efron1994bootstrap}, McNemar's tests \cite{mcnemar1947note}, and Nemenyi tests. This framework provides a rigorous foundation for future Sanskrit NER research.

    \item \textbf{Byte-Level Tokenisation Validation}: We confirmed that byte-level tokenisation (ByT5 \cite{xue2022byt5}) is highly effective for Sanskrit NLP, handling rare compounds, proper names, and sandhi phenomena without vocabulary limitations. This finding has implications for other morphologically rich languages.
\end{enumerate}

\subsection{Methodological Contributions}

Beyond the specific model, this thesis contributes a reusable methodology for low-resource NER:

\begin{itemize}
    \item A transparent, verifiable pipeline for extracting silver NER data from semantic annotation resources (DCS Sembank \cite{hellwig2010dcs})
    \item A principled approach to balancing task-specific performance with preservation of general linguistic capabilities
    \item A three-dimensional evaluation framework (in-domain, linguistic preservation, cross-domain) that can be adapted to other languages and tasks
\end{itemize}

% ============================================================
% 7.2 Key Findings
% ============================================================
\section{Key Findings}

\subsection{Finding 1: Linguistic Regularisation is Highly Effective}

The 15\% linguistic regularisation strategy reduces in-domain F1 by only 0.8 points (0.860 → 0.852) but improves cross-domain PROPN recall by 16 percentage points (57.4\% → 73.4\%). This demonstrates that the regularisation cost is minimal while the generalisation benefit is substantial — a finding with broad implications for task-specific fine-tuning in low-resource settings.

\subsection{Finding 2: Fine-Tuned Models Outperform Zero-Shot LLMs}

The 594M parameter Final NER model significantly outperforms 27B--31B parameter zero-shot large language models (Gemma4 31B \cite{google2025gemma4}: F1 = 0.691; Qwen3.5 27B \cite{qwen2025qwen35}: F1 = 0.672), confirming that task-specific fine-tuning is more effective than scale alone for Sanskrit NER. This challenges the prevailing assumption that larger models are always better and highlights the value of targeted adaptation for low-resource languages.

\subsection{Finding 3: Silver Data Quality Matters More Than Quantity}

The full DCS Sembank \cite{hellwig2010dcs} silver data (12,955 examples from 120 texts) provides substantially greater cross-domain benefit (73.4\% PROPN recall) than limited silver data (2,120 LOC+MISC examples, 57.4\% PROPN recall), despite being automatically extracted. This validates our extraction methodology and suggests that carefully curated silver data can be a valuable resource for low-resource NLP.

\subsection{Finding 4: Trade-offs Are Inherent But Manageable}

The trade-off between in-domain accuracy and cross-domain generalisation is inherent to NER fine-tuning, but our hybrid strategy demonstrates that this trade-off is manageable. The 0.8 F1 point cost of linguistic regularisation is acceptable for most practical applications, particularly those requiring robustness across diverse Sanskrit texts.

% ============================================================
% 7.3 Summary of Key Results
% ============================================================
\section{Summary of Key Results}

\begin{table}[htbp]
    \centering
    \caption{Summary of Key Results: Best NER vs Final NER}
    \label{tab:summary_results}
    \begin{tabular}{lccc}
        \toprule
        \textbf{Metric} & \textbf{Best NER} & \textbf{Final NER} & \textbf{Change} \\
        \midrule
        In-domain F1 (D1) & 0.860 & 0.852 & -0.8 \\
        Cross-domain PROPN Recall (D3) & 57.4\% & 73.4\% & \textbf{+16.0 pp} \\
        Cross-domain Binary F1 (D3) & 0.679 & 0.784 & \textbf{+0.105} \\
        Training Data Size & 133K & 172K & +29\% \\
        Linguistic Regularisation & 0\% & 15\% & — \\
        \bottomrule
    \end{tabular}
\end{table}

% ============================================================
% 7.4 Reproducibility and Broader Impact
% ============================================================
\section{Reproducibility and Broader Impact}

\textbf{Reproducibility}: All code, trained models, and evaluation scripts will be released publicly upon publication. The training scripts (\texttt{train\_final\_ner.py}, \texttt{extract\_full\_sembank\_ner.py}) and all evaluation JSON files are included in the supplementary materials. We hope this facilitates reproducibility and future research in Sanskrit NLP.

\textbf{Broader Impact}: This work contributes to the growing body of research on NLP for low-resource, morphologically rich languages. By demonstrating that hybrid data strategies combining gold annotations, silver data, and linguistic regularisation can achieve competitive performance without requiring massive annotated corpora, we provide a template for other languages with similar resource constraints. The byte-level tokenisation findings may also inform model architecture choices for other Indic languages (Hindi, Bengali, Tamil) where compounding and sandhi phenomena present similar challenges.

% ============================================================
% 7.5 Limitations
% ============================================================
\section{Limitations}

While this thesis makes significant contributions, several limitations should be acknowledged:

\begin{itemize}
    \item \textbf{R\={a}ma/S\={\i}t\={a} Recall Gap}: The DCS Sembank \cite{hellwig2010dcs} NER extraction achieves only 1--2\% recall for major characters due to their dual role as descriptors and entity names.
    \item \textbf{Persistent Class Imbalance}: LOCATION and MISC entities continue to underperform relative to PERSON entities despite oversampling \cite{chawla2002smote}.
    \item \textbf{Genre Skew}: The silver data is heavily skewed toward the Rāmāyaṇa \cite{ramayana} (63\%), limiting diversity across Sanskrit genres.
    \item \textbf{Evaluation Metrics}: Exact match evaluation \cite{tjong2003conll} penalises valid partial matches common in Sanskrit due to sandhi and compounding.
\end{itemize}

These limitations are discussed in detail in Chapter 6 and represent important directions for future improvement.

% ============================================================
% 7.6 Future Work
% ============================================================
\section{Future Work}

\subsection{Short-term Improvements (1--2 years)}

\begin{enumerate}
    \item \textbf{Hybrid Extraction for Dual-Role Entities}: Develop methods to capture entities like Rāma and Sītā that function both as proper names and semantic descriptors without introducing common noun noise.

    \item \textbf{Class-Balanced Loss Functions}: Explore focal loss, class-balanced loss, or synthetic data generation to improve performance on LOCATION and MISC entities.

    \item \textbf{Genre-Diverse Silver Data}: Prioritise extraction from underrepresented genres (medical texts, grammatical works, philosophical śāstras) to improve model robustness.

    \item \textbf{Relaxed Evaluation Metrics}: Develop and validate evaluation metrics (e.g., head-word matching, semantic similarity) more appropriate for Sanskrit philology.
\end{enumerate}

\subsection{Long-term Research Directions (3--5 years)}

\begin{enumerate}
    \item \textbf{Multilingual Indic NER}: Extend the methodology to Hindi, Bengali, Tamil, and other Indic languages, leveraging shared morphological patterns and byte-level tokenisation.

    \item \textbf{Integration with Large Language Models}: Investigate parameter-efficient fine-tuning (LoRA, QLoRA \cite{hu2022lora, dettmers2023qlora}) of 7B--70B parameter models for Sanskrit, combining the strengths of large-scale pretraining with task-specific adaptation.

    \item \textbf{Joint NER and Coreference Resolution}: Develop models that simultaneously perform named entity recognition and coreference resolution, leveraging the CorefUD annotations already present in Mah\={a}n\={a}ma \cite{sarkar2025mahanama}.

    \item \textbf{Domain Adaptation Techniques}: Explore domain adaptation methods (e.g., adversarial training, domain-specific adapters) to improve performance on technical Sanskrit (medical, grammatical, philosophical texts).
\end{enumerate}

\subsection{Dataset and Annotation Efforts}

\begin{enumerate}
    \item \textbf{Expanded Gold Standard}: Create additional gold-standard NER annotations for Vedic \cite{rigveda} Sanskrit, medical texts (Āyurveda \cite{ayurveda}), and grammatical works to reduce domain bias.

    \item \textbf{Sanskrit NER Benchmark}: Establish a public benchmark leaderboard for Sanskrit NER, including multiple test sets (Mahānāma, Pañcatantra \cite{pancatantra_text}, Vedic \cite{rigveda}, technical texts) to drive community progress.

    \item \textbf{Crowdsourced Annotation Platform}: Develop tools and guidelines to enable Sanskrit scholars and students to contribute annotations, following the successful model of Universal Dependencies \cite{nivre2016ud}.
\end{enumerate}

\subsection{Model Architecture Improvements}

\begin{enumerate}
    \item \textbf{Morphology-Aware Models}: Develop models that explicitly incorporate morphological features (case, gender, number) as auxiliary inputs or multi-task objectives.

    \item \textbf{Sandhi-Aware Preprocessing}: Investigate neural sandhi splitting as a preprocessing step to improve entity boundary detection in compounds.

    \item \textbf{Knowledge-Enhanced NER}: Integrate external knowledge sources (Sanskrit WordNet, ontology of characters and locations) to improve entity disambiguation.
\end{enumerate}

% ============================================================
% 7.7 Final Remarks
% ============================================================
\section{Final Remarks}

This thesis has demonstrated that the challenges of Sanskrit Named Entity Recognition — limited gold data, extreme class imbalance, and significant domain diversity — can be effectively addressed through careful methodological design. By combining gold-standard annotations, high-quality silver data, and linguistic regularisation, we have developed a model that achieves competitive in-domain performance while demonstrating strong cross-domain generalisation and preserving linguistic capability.

The work contributes not only a specific model but a broader methodology for low-resource, morphologically rich languages. We hope that the transparent extraction pipeline, the hybrid training strategy, and the comprehensive evaluation framework will serve as a foundation for future research in Sanskrit NLP and related fields.

Ultimately, this thesis advances the goal of \textbf{democratizing access to Sanskrit knowledge} by developing tools that Sanskrit scholars and students can reliably use for large-scale textual analysis. As the field of Sanskrit computational linguistics continues to grow, we look forward to seeing these methods extended, refined, and applied to new texts and new languages.

% ============================================================
% End of Chapter 7
% ============================================================